\documentclass[11pt,a4paper]{article}
\usepackage{fontspec,amsmath,amssymb,amsthm,microtype}
\setmainfont{lmroman10-regular.otf}[BoldFont=lmroman10-bold.otf,ItalicFont=lmroman10-italic.otf,BoldItalicFont=lmroman10-bolditalic.otf]
\usepackage[margin=25mm]{geometry}
\usepackage[unicode,hidelinks]{hyperref}
\setlength{\emergencystretch}{3em}
\setlength{\parskip}{0.4em}
\setlength{\parindent}{0pt}
\allowdisplaybreaks[2]
\raggedbottom
\title{\texorpdfstring{The exact auxiliary velocity estimate in the first-\allowbreak{}order proof}{The exact auxiliary velocity estimate in the first-order proof}}
\author{Omer Deniz Akyildiz\\[0.4em]\normalsize Prepared with AI assistance}
\date{}
\begin{document}\maketitle
Draft mathematical exposition; these notes may contain mistakes.


The source velocity lemma is exposed as an explicit L\ensuremath{{}^{2}} theorem, derived from the synchronous kinetic integral equations and the Gaussian reference law.



The auxiliary exact velocity is bounded using the stationary Gaussian reference and full transformed contraction. This keeps its local velocity radius tied to the current discrepancy rather than repeatedly charging the initial error.


\section{\texorpdfstring{Control the auxiliary velocity by the Gaussian reference and initial phase discrepancy}{Control the auxiliary velocity by the Gaussian reference and initial phase discrepancy}}


Let \(A,R\) be the two continuous kinetic solutions with the same forcing, and put \(e_0=\|P^{-1}(A(0)-R(0))\|_{L^2}\). The synchronous contraction gives \(\|P^{-1}(A(t)-R(t))\|_{L^2}\le e^{-mt/\gamma}e_0\le e_0\) when \(m,t\ge0\). The velocity projection has norm at most the transformed phase norm. The standard Gaussian law of the reference velocity gives \(\|R_V(t)\|_{L^2}=\sqrt p\). The argument applies the triangle inequality to \(A_V(t)=(A_V(t)-R_V(t))+R_V(t)\), substitutes these proved estimates, and obtains \(\|A_V(t)\|_{L^2}\le\sqrt p+e_0\). This is the source auxiliary-velocity lemma, valid without an additional small-step restriction.



The auxiliary and reference exact paths receive the same input, so their full transformed discrepancy contracts from the step’s starting error. Velocity projection has coefficient one in this geometry. Adding the stationary reference velocity radius therefore bounds the auxiliary velocity uniformly throughout the local interval.


\end{document}
